\scanpage{4}
% archive foliation: 3
s'identifiant à l'application qui associe à la droite affine dans $I$ sa
\uline{direction}, qui est une droite de l'espace des translations $V$.
La fibre d'un $\unsure{j}\in\mathcal{P}_{2,2}(I)$, i.e. \unsure{l'ens.} des
\deleted{\unsure{ses}} deux classes \inserted{de cette partition}, est donc
l'espace quotient $I/(D_j=\{0,j\})$.

Le foncteur ``oubli de la structure'' qui va de la catégorie des espaces affines
de dim 2 sur $\mathbb{F}_2$, vers la cat. des ensembles à quatre éléments
\margin{\unsure{on ne prend pour morphismes} les flèches \uline{inversibles}}
\diagram{\begin{tikzcd}
\operatorname{Aff}_2(\mathbb{F}_2) \arrow[r,"\approx"] & (\operatorname{Ens}_4)
\end{tikzcd}}
est une équivalence des catégories.
\inserted{et $=$ un iso, si on définit la notion d'espace affine sur $\mathbb{F}_2$
par la donnée d'une structure sur un $I$ tout seul
(donnée d'un sous-groupe de $\unsure{\mathcal{G}_I|\mathfrak{S}_I}$ \unsure{gr de tr.})}

Je dis que le foncteur
\diagram{\begin{tikzcd}
(\operatorname{Ens}_4) \arrow[r] & (\operatorname{Gr}) \\
I \arrow[r,mapsto] & \unsure{\mathcal{G}_I|\mathfrak{S}_I}
\end{tikzcd}}
est pleinement fidèle, ou encore que pour $\forall I$,
$\unsure{\mathcal{G}_I|\mathfrak{S}_I}$ a un centre trivial et tout automorphisme de
$\unsure{\mathcal{G}_I|\mathfrak{S}_I}$ est intérieur.
\deleted{\unsure{Pour}}\inserted{\unsure{Il}}
\unsure{suffirait} bien sûr de l'énoncer pour $\unsure{\mathcal{G}_4|\mathfrak{S}_4}$ !).

On va définir pour ceci un foncteur quasi-inverse
\diagram{\begin{tikzcd}
\deleted{\unsure{(\operatorname{Ens})}}(\operatorname{Gr})'
\arrow[r] & (\operatorname{Ens}_4)
\end{tikzcd}}
sur l'image essentielle $(\operatorname{Gr})'$ des groupes isomorphes à
$\unsure{\mathcal{G}_4|\mathfrak{S}_4}$.
Si on a un tel groupe $G$, on considère l'ens
\deleted{\unsure{des}}\deleted{\unsure{sous-groupes de}}
\inserted{3-sous-gr.} de Sylow, il y en a exactement quatre !
Pour \unsure{préciser} ce point et réciproquement
